%% =====================================================================
%%  T-22-02  Borrowed Warmth
%%  -------------------------------------------------------------------
%%  One idea per file. Core is mandatory. Slots in canonical order.
%%  Max 700 words. No layout commands. No filler. New source material
%%  for this entry goes in sources/<BOOK>/T-22-02.tex -- never by
%%  rewriting this file (docs/RULES.md R0.1).
%% =====================================================================
%% @kind: topic
%% @id: T-22-02
%% @title: Borrowed Warmth
%% @subtitle: Attaching yourself to what is already liked
%% @chapter: c22
%% @order: 20
%% @status: stable
%% @sources: B4
%% @updated: 2026-09-19

\topic{T-22-02}{Borrowed Warmth}{Attaching yourself to what is already liked}

\begin{core}
Regard transfers by contact. A person, product or proposal that appears alongside something already valued is valued more highly, and the transfer happens without any argument being made. The same mechanism runs in reverse, which is why people distance themselves from bad news they had no part in.
\end{core}
\begin{mechanism}
The transfer works because association is one of the ways the mind economises on evaluation. Judging something on its merits requires information; judging it by what it is next to requires only attention. Where information is unavailable --- which is most of the time --- the second method is used and is experienced as the first.

This makes proximity a controllable variable rather than an accident. The messenger can be chosen for their warmth, the setting for its quality, the photograph for its company, and the announcement for its timing alongside good news. It also explains a consistent asymmetry: people claim connection with success and disclaim connection with failure, using first-person language for one and distant language for the other. That behaviour is the same mechanism observed from the other side, and it is a reliable diagnostic in itself.
\end{mechanism}
\begin{tells}
  \item A proposal arrives attached to a name, a cause or an event you already approve of, and the connection is emphasised.
  \item The messenger was chosen for how you feel about them rather than for what they know.
  \item Good news and the request arrive together.
  \item Someone uses first-person language about a success they did not produce, and distant language about a failure they did.
  \item You like the proposal and cannot separate it from its company.
\end{tells}
\begin{moves}
  \item Choose the messenger for the counterpart's regard, not for relevance.
  \item Time the request to follow good news rather than precede it.
  \item Place the proposal alongside something already valued --- a shared cause, a respected name, a familiar setting.
  \item Claim successes in the first person and describe failures distantly; both are read as evidence about you.
\end{moves}
\begin{counters}
  \item Detach the proposal and evaluate it alone. Ask what it would look like presented by someone you distrust, in a plain room, on an ordinary day.
  \item Notice the messenger's selection. A person chosen for their warmth was chosen for a reason, and the reason is not their expertise.
  \item Watch your own language about successes and failures. The asymmetry is the tell, in yourself as much as in others.
  \item Separate the announcement from the timing. Ask what else happened that day.
\end{counters}
\begin{cost}
Borrowed warmth is rented, not owned, and the rent comes due when the borrowed thing fails --- at which point the discredit transfers back with the same mechanism that carried the credit. Building a position entirely out of association also leaves nothing standing when the associations are removed, which is a common failure in careers built on proximity rather than on output.
\end{cost}

\sources{B4}
\seesources
\seealso{T-22-01, T-04-03, T-09-01, T-23-02}
